\section{Implementation Code Listings and Deployment Artifacts}

This appendix presents key implementation excerpts, database schemas, and operational commands central to the security, privacy, and reliability guarantees of SIGS-RS.

\subsection{Fail-Secure Privacy Anonymizer Implementation}
Listing~\ref{lst:anonymizer_impl} illustrates the core anonymization method from \texttt{src/sigs\_rs/middleware/anonymizer.py}, showcasing deep copying, recursive PII stripping, post-sanitization residual assertions, and fail-secure exception handling (\texttt{REQ-PA-1}, \texttt{REQ-PA-4}).

\begin{figure}[H]
\centering
\begin{minipage}{0.95\linewidth}
\begin{verbatim}
class PrivacyAnonymizer:
    def anonymize(self, packet: RawTelemetryPacket) -> SanitizedTelemetryPacket:
        sensor_id = getattr(packet, "sensor_id", "<unknown>")
        try:
            if not isinstance(packet, RawTelemetryPacket):
                raise TypeError("expected a verified RawTelemetryPacket")

            # Work from a detached copy: failed mask never alters raw packet
            raw = deepcopy(packet.model_dump())
            self._strip_pii(raw, self._pii_fields)
            self._assert_no_pii(raw, self._pii_fields)

            # Whitelist only fields present on safe output type
            sanitized_fields = {
                k: raw[k] for k in SanitizedTelemetryPacket.model_fields if k in raw
            }
            sanitized_fields["ingested_at"] = datetime.now(timezone.utc)
            return SanitizedTelemetryPacket(**sanitized_fields)
        except AnonymizerError:
            raise
        except Exception as exc:
            raise AnonymizerError(
                f"Masking failed for sensor '{sensor_id}': {exc}"
            ) from exc

    @staticmethod
    def _strip_pii(data: Any, pii_fields: list[str]) -> None:
        if isinstance(data, dict):
            targets_lower = {f.lower() for f in pii_fields}
            keys_to_remove = [k for k in list(data.keys()) if k.lower() in targets_lower]
            for field in keys_to_remove:
                data.pop(field, None)
            for value in data.values():
                PrivacyAnonymizer._strip_pii(value, pii_fields)
        elif isinstance(data, list):
            for value in data:
                PrivacyAnonymizer._strip_pii(value, pii_fields)
\end{verbatim}
\end{minipage}
\caption{Fail-Secure Privacy Anonymizer Logic (\texttt{anonymizer.py})}
\label{lst:anonymizer_impl}
\end{figure}

\subsection{Database Schema and Immutable Audit Trigger}
Listing~\ref{lst:database_schema_sql} illustrates the PostgreSQL/TimescaleDB schema definition in \texttt{schema.sql}, demonstrating daily chunk hypertable creation for sanitized readings and PL/pgSQL trigger enforcement preventing \texttt{UPDATE} or \texttt{DELETE} tampering on the audit log (\texttt{REQ-PA-3}, \texttt{REQ-SE-4}).

\begin{figure}[H]
\centering
\begin{minipage}{0.95\linewidth}
\begin{verbatim}
-- Create TimescaleDB hypertable for sanitized readings
CREATE TABLE IF NOT EXISTS telemetry.sanitized_readings (
    id              BIGSERIAL,
    ts              TIMESTAMPTZ      NOT NULL,
    ingested_at     TIMESTAMPTZ      NOT NULL DEFAULT NOW(),
    sensor_id       VARCHAR(64)      NOT NULL,
    resource_type   VARCHAR(32)      NOT NULL,
    value           DOUBLE PRECISION NOT NULL,
    unit            VARCHAR(32)      NOT NULL,
    metadata        JSONB,
    PRIMARY KEY (id, ts)
);

SELECT create_hypertable(
    'telemetry.sanitized_readings', 'ts',
    if_not_exists => TRUE, chunk_time_interval => INTERVAL '1 day'
);

-- Immutable security audit log table
CREATE TABLE IF NOT EXISTS audit.security_log (
    id          BIGSERIAL        PRIMARY KEY,
    ts          TIMESTAMPTZ      NOT NULL DEFAULT NOW(),
    event_type  VARCHAR(64)      NOT NULL,
    severity    VARCHAR(16)      NOT NULL DEFAULT 'INFO',
    sensor_id   VARCHAR(64),
    user_id     VARCHAR(128),
    detail      JSONB            NOT NULL DEFAULT '{}',
    ip_address  VARCHAR(45)
);

-- PL/pgSQL trigger blocking modification or deletion (SRS REQ-SE-4)
CREATE OR REPLACE FUNCTION audit.prevent_tampering()
RETURNS TRIGGER AS $$
BEGIN
    RAISE EXCEPTION 'audit.security_log is immutable: UPDATE and DELETE prohibited';
END;
$$ LANGUAGE plpgsql;

CREATE TRIGGER trg_audit_immutable
    BEFORE UPDATE OR DELETE ON audit.security_log
    FOR EACH ROW
    EXECUTE FUNCTION audit.prevent_tampering();
\end{verbatim}
\end{minipage}
\caption{PostgreSQL TimescaleDB Hypertable and Immutable Audit Trigger Schema}
\label{lst:database_schema_sql}
\end{figure}

\subsection{Selected System Commands and Run Instructions}
\begin{itemize}
    \item \textbf{Execute Hermetic Live Demonstration:}
\begin{verbatim}
.venv/bin/python -m demo.run_live_demo
\end{verbatim}
    \item \textbf{Run Complete Automated Verification Suite (580 Tests):}
\begin{verbatim}
.venv/bin/pytest tests/unit/ tests/integration/ tests/benchmarks/ -v
\end{verbatim}
    \item \textbf{Execute Regression Invariant Verification (Phase 16 Hardening):}
\begin{verbatim}
.venv/bin/pytest tests/unit/test_phase16_audit_hardening.py -v
\end{verbatim}
    \item \textbf{Start Standalone Local Demo Server (Port 8000):}
\begin{verbatim}
.venv/bin/python -m demo.demo_server
\end{verbatim}
    \item \textbf{Execute Local Performance Microbenchmarks (Stages 0--2):}
\begin{verbatim}
.venv/bin/python -m benchmarks.cli --suite all --stages 0,1,2
\end{verbatim}
\end{itemize}
